% Part: sets-functions-relations
% Chapter: sets
% Section: basics

\documentclass[../../../include/open-logic-section]{subfiles}

\begin{document}

\olfileid{sfr}{set}{bas}
\olsection{మూలకాధారిత సమానత్వం}

వస్తువుల సమూహాన్ని {ఒకే} వస్తువుగా పరిగణించినప్పుడు దానిని
\emph{సమితి} అంటాం. ఆ సమితిని ఏర్పరచే వస్తువులను దాని
\emph{మూలకాలు} లేదా \emph{సభ్యాలు} అంటాం. $x$ అనేది $A$ అనే
సమితిలో !!a{element} అయితే $x \in A$ అని రాస్తాం; కాకపోతే
$x \notin A$ అని రాస్తాం. !!{element}s ఏవీ లేని సమితిని
\emph{శూన్య} సమితి అని పిలిచి, ``$\emptyset$''తో సూచిస్తాం.

\begin{explain}
సమితిని ఏ విధంగా \emph{నిర్దేశించాం}, దాని
\tecase{acc}{!!{element}s} ఏ \emph{క్రమంలో} పేర్కొన్నాం, లేదా
ఆ \tecase{acc}{!!{element}s} \emph{ఎన్నిసార్లు} లెక్కించాం అన్నది ముఖ్యం కాదు.
దాని !!{element}s ఏమిటన్నదే ముఖ్యం. ఈ విషయాన్ని కింది సూత్రంగా
వ్యక్తం చేస్తాం.
\end{explain}

\begin{defn}[మూలకాధారిత సమానత్వం]
  $A$, $B$ సమితులైతే, $A = B$ కావడానికి అవసరమైన మరియు సరిపడే
  షరతు: $A$లోని ప్రతి !!{element}, $B$లో కూడా !!a{element}
  కావాలి; తిరుగుదిశలోనూ ఇదే వర్తించాలి.
\end{defn}

మూలకాధారిత సమానత్వం కొన్ని సంకేతాలను వాడటానికి ఆధారమిస్తుంది.
సాధారణంగా $a_{1}$, \dots, $a_{n}$ అనే వస్తువులు ఉన్నప్పుడు,
$\{a_{1}, \dots, a_{n}\}$ అనేది $a_1, \ldots, a_n$ మాత్రమే
!!{element}sగా గల \emph{ఆ ఏకైక} సమితి. ఇక్కడ
``\emph{ఆ ఏకైక}'' అన్న మాటను నొక్కిచెబుతున్నాం: మూలకాధారిత
సమానత్వం ప్రకారం అటువంటి సమితి \emph{ఒక్కటే} ఉండగలదు.
అదే సూత్రం ప్రకారం కిందిదీ నిజమే:
  \[
    \{a, a, b\} = \{a, b\} = \{b,a\}.
  \]
సమితుల విషయంలో వాటి \tecase{gen}{!!{element}s} క్రమాన్ని గానీ,
ఒక మూలకాన్ని ఎన్నిసార్లు పేర్కొన్నామన్న విషయాన్ని గానీ పరిగణించనవసరం
లేదని ఇది స్పష్టం చేస్తుంది.

\begin{tagblock}{novice}
\begin{ex}
కొన్ని వస్తువులు ఉన్నప్పుడు వాటిని ఒక సమితిగా సమీకరించవచ్చు.
ఉదాహరణకు, రిచర్డ్ సహోదరుల సమితిలో ఒక వ్యక్తి ఉన్నారు;
దానిని $S=\{\textrm{Ruth}\}$ అని రాయవచ్చు.
$4$ కంటే చిన్న ధన పూర్ణ సంఖ్యల సమితి $\{1, 2, 3\}$.
దానిని $\{3, 2, 1\}$ అని, లేదా $\{1, 2, 1, 2, 3\}$ అని కూడా
రాయవచ్చు. మూలకాధారిత సమానత్వం ప్రకారం ఇవన్నీ ఒకే సమితి.
ఎందుకంటే $\{1, 2, 3\}$లోని ప్రతి !!{element}, $\{3, 2, 1\}$లో
(అలాగే $\{1, 2, 1, 2, 3\}$లో) !!a{element}; తిరుగుదిశలోనూ
ఇదే వర్తిస్తుంది.
\end{ex}
\end{tagblock}

తరచుగా ఒక సమితిలోని !!{element}s అన్నింటికీ ఉమ్మడిగా ఉన్న
ధర్మం ద్వారా ఆ సమితిని నిర్దేశిస్తాం. ఇందుకు
$\Setabs{x}{\phi(x)}$ అనే సంక్షిప్త సంకేతాన్ని వాడతాం.
ఇందులో $\phi(x)$ అనేది ఆ సమితిలోని \tecase{loc}{!!{element}s}
లెక్కించబడాలంటే $x$కు ఉండవలసిన ధర్మాన్ని సూచిస్తుంది.

\begin{tagblock}{novice}
\begin{ex}
మన ఉదాహరణలో $S$ను ఇలా కూడా నిర్దేశించవచ్చు:
\[
S = \Setabs{x}{x \text{ రిచర్డ్‌కు సహోదర సంబంధం గల వ్యక్తి}}.
\]
\end{ex}
\end{tagblock}

\begin{tagblock}{math}
\begin{ex}
ఒక సంఖ్య తన నిజ భాజకాల మొత్తానికి సమానమైతే, అప్పుడే దానిని
\emph{పరిపూర్ణ} సంఖ్య అంటాం. నిజ భాజకాలు అంటే ఆ సంఖ్యను
నిశ్శేషంగా భాగించే, కానీ ఆ సంఖ్యకు సమానం కాని సంఖ్యలు.
ఉదాహరణకు $6$ పరిపూర్ణ సంఖ్య: దాని నిజ భాజకాలు $1$, $2$, $3$;
అలాగే $6 = 1 + 2 + 3$. వాస్తవానికి $10$ కంటే చిన్న ధన పూర్ణ
సంఖ్యలలో పరిపూర్ణమైనది $6$ ఒక్కటే. కాబట్టి మూలకాధారిత సమానత్వంతో
ఇలా చెప్పవచ్చు:
\[
	\{6\} = \Setabs{x}{x\text{ పరిపూర్ణం మరియు }0 \leq x \leq 10}
\]
కుడివైపు సంకేతాన్ని ``$x$ పరిపూర్ణమై, $0 \leq x \leq 10$ అయ్యే
$x$ల సమితి'' అని చదువుతాం. సమితులను ఏ విధంగా నిర్దేశించామన్నది
వాటి సమానత్వానికి ముఖ్యం కాదని ఇక్కడి సమానత నిర్ధారిస్తుంది.
మరింత సాధారణంగా, $\phi(x)$ అయ్యే $x$ల సమితి ఒక్కటే ఉంటుందని
మూలకాధారిత సమానత్వం హామీ ఇస్తుంది.
అందువల్ల $\Setabs{x}{\phi(x)}$ను $\phi(x)$ అయ్యే $x$ల
\emph{ఆ ఏకైక} సమితి అని పిలవడానికి ఈ సూత్రం ఆధారమిస్తుంది.
\end{ex}
\end{tagblock}

రెండు సమితులు ఒకటేనని చూపే పద్ధతిని మూలకాధారిత సమానత్వం ఇస్తుంది:
$A = B$ అని చూపాలంటే, $x \in A$ అయినప్పుడల్లా $x \in B$
అవుతుందని, $y \in B$ అయినప్పుడల్లా $y \in A$ అవుతుందని చూపాలి.

\begin{prob}
శూన్య సమితి మహా అయితే ఒక్కటే ఉండగలదని నిరూపించండి. అంటే,
!!{element}s ఏవీ లేని సమితులు $A$, $B$ అయితే $A = B$ అని చూపండి.
\end{prob}

\end{document}
